\documentclass[10pt]{article}
\usepackage{amsmath,amssymb}
\usepackage[margin=1in]{geometry}
\begin{document}
\section*{Duality}
For every primal feasible $x \in \mathbb{R}^n$ and dual feasible $y$, weak
duality gives $c^\top x \ge b^\top y$, and the gap $\sqrt{\sum_i (x_i - y_i)^2}$
closes at an optimum. The tableau
\[
  T = \left(\begin{array}{ccc|c} 1 & 0 & -2 & 4 \\ 0 & 1 & 3 & 5 \\ 0 & 0 & 1 & 6 \end{array}\right),
  \qquad
  \Bigg\{ \sum_{j=1}^{n} a_{ij} x_j \le b_i \Bigg\}
\]
is optimal when every reduced cost is non-negative. Complementary slackness
requires $x_j \bigl( c_j - \sum_i a_{ij} y_i \bigr) = 0$ for every $j$, and the
ratio test picks $\min_i \{ b_i / a_{ij} : a_{ij} > 0 \}$. Integrals such as
\[
  \int_0^\infty e^{-\lambda t}\, dt = \frac{1}{\lambda},
  \qquad
  \left\lVert \frac{\partial f}{\partial x} \right\rVert \le \sqrt{\frac{L}{\mu}}
\]
appear in the convergence analysis. Each line of ordinary text in this
paragraph should select as a tight band around its own glyphs, without the
highlight of one line covering the letters of the line above or below it,
even though the extension font draws delimiters three ems tall.
\end{document}
